\documentclass[11pt,a4paper]{article}
\usepackage[margin=2.5cm]{geometry}
\usepackage{amsmath,amssymb}
\setlength{\emergencystretch}{4em}

% --- notation shared with traceable_energetics.tex ----------------------------------
\newcommand{\vt}{\vartheta}
\newcommand{\code}[1]{\texttt{#1}}
\newcommand{\pdc}[3]{\left.\frac{\partial #1}{\partial #2}\right|_{#3}}
\newcommand{\Iz}{\mathcal{I}_z}
% this note
\newcommand{\Mtot}{M_{\rm tot}}
\newcommand{\Ptop}{P_{\rm top}}
\newcommand{\zs}{z_{*}}
\newcommand{\zh}{\hat z}

\title{Reference states for available potential energy\\ in AN2D and FC2D: algorithms}
\author{Notes on the \code{TraceableEnergetics} code}
\date{}

\begin{document}
\maketitle
\tableofcontents

\bigskip
\noindent\textbf{Scope.}
This note specifies how to compute the \emph{reference state} of a snapshot of AN2D or FC2D, precisely enough to be
implemented. The available energy density, its budget and the conditions that the reference state must satisfy are derived
in the energy-budget notes \code{energy.tex} (``Energy budgets for compressible and anelastic fluids'', ``the energy notes''
below). Section~\ref{sec:theory} summarises what they require, and \S\ref{sec:outputs} lists what this note delivers to them.
Notation and code names follow \code{traceable\_energetics.tex} (``the main note''): $\vt$ is the entropic variable
(\code{consvar}), $q$ the solute fraction, $h(P,\vt,q)$ the specific enthalpy with $\partial h/\partial P|_{\vt,q}=v$,
$\Phi=g(z-z_0)$, and the box is $[0,L_x]\times[0,L_z]$ with $M_x\times M_z$ interior cells of size $\Delta x\times\Delta z$.

\paragraph{Correspondence with the energy notes.} The energy notes write the pressure $p$ and $h(\vt,q,p)$. Here, as in the
main note, $P$ is the pressure and $p=(P-P_R)/\rho_R$ the reduced pressure of AN2D. Their static energy is
\begin{equation}
  \Sigma(\vt,q,z,P)=\begin{cases}
    h(P,\vt,q)+\Phi(z) & \text{FC2D},\\[2pt]
    h(P_R(z),\vt,q)+\dfrac{P-P_R(z)}{\rho_R(z)}+\Phi(z) & \text{AN2D}
  \end{cases}
  \label{eq:Sigma}
\end{equation}
(they take $z_0=0$; the available energy does not depend on $z_0$). Reference-state quantities carry a subscript $*$, as
there: the reference pressure $P_*(z,t)$, the reference height $\zs(\vt,q,t)$ of a parcel, the reference specific volume
$v_*(z,t)$, the reference density $\rho_*=1/v_*$ (FC2D) and $N_*^2(\vt,q,t)$. In AN2D the reduced reference pressure is
$\pi_*=(P_*-P_R)/\rho_R$.

\medskip
Two families of reference state are covered:
\begin{enumerate}
  \item the \textbf{Lorenz reference state}: the state of minimum potential energy reachable by an adiabatic
        rearrangement of the parcels (\S\ref{sec:lorenz}--\ref{sec:histogram}). Four solvers are specified: sorting,
        greedy filling with repair, exact rearrangement, and histogram (PDF) reduction;
  \item the \textbf{horizontal-mean reference state}: a profile that is not a rearrangement of the fluid
        (\S\ref{sec:mean}).
\end{enumerate}
In both, parcels are given their reference height by the level-of-neutral-buoyancy (LNB) equation against the reference
profile (\S\ref{sec:lnb-sec}). Section~\ref{sec:decisions} records the design decisions, including an accepted
discretisation mismatch in FC2D, and \S\ref{sec:tests} lists the tests.

% =====================================================================================
\section{What the reference state must satisfy}\label{sec:theory}
The energy notes define the available energy (AE) of a parcel at height $z$ and pressure $P$,
\begin{equation}
  A=\Sigma(\vt,q,z,P)-\Sigma\left(\vt,q,\zs,P_*(\zs,t)\right)-\left(P-P_*(z,t)\right)\Sigma_P(\vt,q,z,P),
  \label{eq:A}
\end{equation}
with $\Sigma_P=v$ (FC2D) or $1/\rho_R$ (AN2D). It splits as $A=\Pi_1+\Pi_2$ with
$\Pi_2=\Sigma(\vt,q,z,P_*(z,t))-\Sigma(\vt,q,\zs,P_*(\zs,t))$, and $\Pi_1=0$ in AN2D. The reference profile defines $v_*$
by
\begin{equation}
  v_*\,\frac{\partial P_*}{\partial z}=-g\quad\text{(FC2D)},\qquad
  \frac{\partial\pi_*}{\partial z}=g\left(\rho_R v_*-1\right)\quad\text{(AN2D)}.
  \label{eq:vstar}
\end{equation}
The conditions on $P_*$ and $\zs$ (energy notes, ``Reference state'' and the appendix ``Variational origin of the reference
state conditions'') are:
\begin{enumerate}
  \item \textbf{Level of neutral buoyancy (LNB).} $\partial_z\Sigma(\vt,q,z,P_*(z,t))=0$ at $z=\zs$, i.e.\
        $v(P_*(\zs,t),\vt,q)=v_*(\zs,t)$ (FC2D) or $v(P_R(\zs),\vt,q)=v_*(\zs,t)$ (AN2D).
  \item \textbf{Maximal APE.} $N_*^2=\partial_z^2\Sigma(\vt,q,z,P_*(z,t))>0$ at $z=\zs$.
  \item \textbf{Volume constraint.} For every $f(z)$,
        \begin{equation}
          \int f(\zs)\,\Sigma_P\left(\vt,q,\zs,P_*(\zs,t)\right)\mathrm{d}m=\int f(z)\,\mathrm{d}V ,
          \label{eq:volume}
        \end{equation}
        i.e.\ the reference state fills the box level by level. In AN2D $\Sigma_P=1/\rho_R(\zs)$, so the reference state has
        the mass per unit height $L_x\rho_R(z)$ of the model.
\end{enumerate}
These are the stationarity conditions of $\int A\,\mathrm{d}m$ with respect to the reference state. At fixed $P_*$, $A$
depends on $\zs$ only through $-\Sigma(\vt,q,\zs,P_*(\zs,t))$, so (1) and (2) are the first- and second-order conditions for
$\zs$ to maximise the parcel's $A$. Condition (3) is stationarity with respect to $P_*$, which is the Lagrange multiplier of the
volume constraint (exactly so in AN2D; in FC2D $\int A\,\mathrm{d}m$ is a minimum with respect to $P_*$). When (3) holds,
$\int A\,\mathrm{d}m=E_P-E_*$, the potential energy of the actual state minus that of the rearranged state (Lorenz). With the
choice of $\zs$ below, the rearranged state is then the one of minimum energy (\S\ref{sec:dual}).

\paragraph{Choice of $\zs$.} This note takes $\zs$ to be the \emph{global} minimiser of $z\mapsto\Sigma(\vt,q,z,P_*(z,t))$
on $[0,L_z]$, the height that maximises the parcel's $A$. It is a function of $(\vt,q,t)$ only, as in the energy notes, and it
satisfies (1)--(2) whenever it is interior. Section~\ref{sec:lnb} gives the algorithm and compares it with the rule of
Tailleux (2013).

\paragraph{What the budgets need.} Beyond the model state, the AE budgets of the energy notes need $\zs$ for each parcel,
$P_*(z,t)$ and $v_*(z,t)$ at the model rows and at $\zs$, and $\partial_tP_*(z,t)$ (\S\ref{sec:dtpstar}). The terms in
$\partial_t\zs$ drop out of the budgets by the LNB condition, so $\zs$ need not be followed in time. This holds only if $\zs$
minimises $\Sigma(\vt,q,z,P_*(z,t))$ for the same profile $P_*$ that is used to evaluate $A$ (\S\ref{sec:profile}).

\paragraph{The discrete problem.} Section~\ref{sec:lorenz} poses the Lorenz state as an optimal rearrangement of the parcels between levels. Its placement is
the discrete $\zs$, its slot constraints are the discrete volume constraint, and its dual potentials are the discrete $P_*$
(\S\ref{sec:dual}). The rearrangement problem and its dual are the discrete form of the variational principle above.

% =====================================================================================
\section{Why the problem is not just a sort}\label{sec:why}
For a single-component fluid, or a Boussinesq fluid with an equation of state linear in its tracers, the Lorenz state is
obtained by sorting the parcels by a single conserved quantity (Winters et al.\ 1995). The fluid here is binary and its
equation of state (main note, Appendix~A) is
\begin{equation}
  v = v_0\left[1+\alpha_T\left(1+\gamma(P-P_0)\right)(\vt-T_0)-\alpha_q(q-q_0)-\alpha_P(P-P_0)+\tfrac12\alpha_{TT}(\vt-T_0)^2\right].
  \label{eq:eos}
\end{equation}
\paragraph{When sorting is exact.} Suppose $v$ is additively separable,
\begin{equation}
  v(P,\vt,q)=f(\vt,q)+w(P).
  \label{eq:separable}
\end{equation}
Then $h(P,\vt,q)=h(P_0,\vt,q)+f(\vt,q)\,(P-P_0)+\int_{P_0}^{P}w$, and the cost of placing parcel $i$ at pressure $P$ is
$a_i+f_i P+(\text{a function of }P\text{ only})$. Only the term $\sum f_iP$ depends on the arrangement, and by the
rearrangement inequality it is minimised by placing the largest $f$ (lightest parcel) at the lowest pressure (the top). The
same holds with unequal masses and split parcels: with the parcels ordered by $f$ and the levels by pressure, the costs have
the Monge property, for which filling the levels in order (the north-west corner rule) solves the rearrangement problem of
\S\ref{sec:rearrangement} exactly (Hoffman 1963). In FC2D the pressures of the levels also depend on $\Ptop$
(\S\ref{sec:fc-slots}), but the order does not, so the same placement is optimal for every $\Ptop$. So \textbf{the sort by $f$ is
the exact minimiser whenever \eqref{eq:separable} holds}. For \eqref{eq:eos} this is the case iff $\gamma=0$:
compressibility ($\alpha_P$), haline contraction ($\alpha_q$) and cabbeling ($\alpha_{TT}$) do not break it;
\textbf{thermobaricity ($\gamma\neq0$) does}. The experiments use $\gamma=10^{-7}\,$Pa$^{-1}$, so the sort is an
approximation. The relative size of the coupling term across the box is $\gamma\,\Delta P\approx\gamma\rho g L_z$
($\approx5\times10^{-4}$ for $L_z=0.5\,$m), so the sort should be very close to exact: only parcels whose specific volumes are
nearly tied can be misordered. This must be checked (\S\ref{sec:sort}, \S\ref{sec:exact}), not assumed.

\paragraph{The general problem.} Without \eqref{eq:separable}, the order in which parcels stack depends on the pressure at
which they are compared. Finding the minimum is then a linear programme, known in the literature as the \emph{linear assignment} or \emph{transportation} problem
(Hieronymus \& Nycander 2015), solvable exactly in polynomial time (\S\ref{sec:exact}). Heuristics that fill the column
level by level can return a ``reference'' state with \emph{more} energy than the actual state, i.e.\ a negative APE (Harris \&
Tailleux 2018 found this in 30--40\% of atmospheric soundings for top-down filling), so every heuristic must be followed by
an energy check (\S\ref{sec:greedy}).

% =====================================================================================
\section{The discrete minimisation problem}\label{sec:lorenz}
\subsection{Parcels}
Each interior cell $(i,j)$ is a parcel, indexed by a single $n=1,\dots,N$, $N=M_xM_z$, with mass and conserved properties
\begin{equation}
  m_n=\begin{cases}\rho_{R,j}\,\Delta x\Delta z & \text{AN2D},\\ m^\rho_{ij}\,\Delta x\Delta z & \text{FC2D},\end{cases}\qquad
  \vt_n=\frac{m^\vt_{ij}}{\rho},\quad q_n=\frac{m^q_{ij}}{\rho},
\end{equation}
with $\rho=\rho_{R,j}$ (AN2D) or $m^\rho_{ij}$ (FC2D). The total mass is $\Mtot=\sum_nm_n$. A rearrangement keeps each
parcel's $(m_n,\vt_n,q_n)$, which is what adiabatic, isohaline, mass-conserving motion does. The reference state is at rest
and fills the box. Its pressure $P_*(z,t)$, hence $v_*(z,t)$, is horizontally uniform, but $(\vt,q)$ need not be uniform
along a level (density compensation; energy notes).

\subsection{Slots and the rearrangement problem}\label{sec:rearrangement}
In both models the reference state is a stack of horizontal layers spanning $L_x$. Its natural vertical coordinate is
\emph{mass}, because the pressure (FC2D) and the height (AN2D) of a layer are known once the mass on either side of it is
known, whatever parcels occupy it. Divide the column into $K$ \emph{slots} $k=1,\dots,K$ of mass $M_k$
($\sum_kM_k=\Mtot$), each with a per-unit-mass cost $c_{nk}$ of placing parcel $n$ there. A rearrangement is a
placement $x_{nk}\ge0$, the mass of parcel $n$ placed in slot $k$:
\begin{equation}
  \min_{x\ge0}\ \sum_{n,k}x_{nk}\,c_{nk}\qquad\text{subject to}\qquad \sum_kx_{nk}=m_n\ \ \forall n,\qquad \sum_nx_{nk}=M_k\ \ \forall k .
  \label{eq:tp}
\end{equation}
This linear programme is known as the transportation problem. When all $m_n$ and $M_k$ are equal and $K=N$ it reduces to the
linear assignment problem of Hieronymus \& Nycander (2015). With $K\ll N$ shared levels it is the minimum-cost flow problem
of Su \& Ingersoll (2016). Parcels may be split between slots, which is needed because parcel masses differ.
The placement is the discrete $\zs$: $x_{nk}>0$ places (part of) parcel $n$ at the level of slot $k$.
The models differ only in their slots and costs.

\subsection{AN2D: slots are the model rows}\label{sec:an-slots}
By the volume constraint \eqref{eq:volume} with $\Sigma_P=1/\rho_R$, the reference state has the mass per unit height
$L_x\rho_R(z)$ of the model. These are also the states that anelastic flow can reach: it satisfies $\nabla\cdot(\rho_R\mathbf u)=0$,
so its flow map preserves the measure $\rho_R\,\mathrm{d}A$. So the slots are the model rows, and the thermodynamics is
evaluated at $P_R(z)$, as in the model. With one slot per row,
\begin{equation}
  K=M_z,\qquad M_k=M_x\,\rho_{R,k}\,\Delta x\Delta z,\qquad c_{nk}=h(P_{R,k},\vt_n,q_n)+g\,(z_k-z_0).
  \label{eq:an-cost}
\end{equation}
The slot resolution is set by the user: with $s$ slots per row ($K=sM_z$), row $k$ is split into $s$ slots of mass $M_k/s$
and height $\Delta z/s$, with the costs at their centres and $P_R$ from the interpolant of \S\ref{sec:profile}. The statements
below on exact discrete consistency hold for $s=1$. For $s>1$ the actual state is no longer a placement with exactly the model's
energy, and they hold to $O(\Delta z^2)$, as in FC2D.
The reference pressure $P_*=P_R+\rho_R\pi_*$ is not $P_R$: it is the multiplier of the slot constraints (\S\ref{sec:dual}).
The potential energy of a placement, with $\mathcal U=\rho_R(h(P_R)+\Phi)-P_R$ of the main note, is
\begin{equation}
  E[x]=\sum_{n,k}x_{nk}\,c_{nk}-\sum_kM_x\,P_{R,k}\,\Delta x\Delta z ,
\end{equation}
where the last term does not depend on the placement. The actual state is the feasible placement ``each parcel in its own row'', and its
energy $E_P$ is exactly the discrete $\int\mathcal U$ of the model. Hence \textbf{the exact solution gives $E_P\ge E_*$ to
round-off, with equality when the actual state is itself a minimiser}: a horizontally uniform state that is stable to finite
exchanges of parcels (with \eqref{eq:separable}, any horizontally uniform, statically stable state). The reference state is
consistent with the model's discrete energy.

\paragraph{Placement height of a parcel.} A parcel split over several rows has the mass-weighted height
$\zh_n=\sum_kx_{nk}z_k/m_n$. Inside a row the parcels are not ordered by the placement. If a sub-row height is wanted
(the analogue of the Winters et al.\ height), order the parcels within row $k$ by the sort key of \S\ref{sec:sort} and place
them by cumulative mass: $\zh_n=z_k-\tfrac12\Delta z+(m_{<n}+\tfrac12x_{nk})/(L_x\rho_{R,k})$, where $m_{<n}$ is the mass
in the row below parcel $n$. The $\zs$ used in $A$ is the LNB height of \S\ref{sec:lnb}; $\zh$ is used to check it.

\paragraph{Reference profile.} The reference state has no single $(\vt,q)$ per row. Its reduced pressure on the slot centres
is given by the duals of the slot constraints, $\pi_{*,k}=-\psi_k$ (\S\ref{sec:dual}). This is the exact multiplier, for which
$\int A\,\mathrm{d}m=E_P-E_*$ exactly ($s=1$). A solver that returns no duals (greedy+repair, \S\ref{sec:greedy}) falls back
on the slot-mean specific volume
\begin{equation}
  v_{*,k}=\frac{1}{M_k}\sum_nx_{nk}\,v(P_{R,k},\vt_n,q_n),
\end{equation}
with $\pi_*$ hydrostatic with it between the slot centres by the trapezoidal rule. For $s=1$ this is the face average $\Iz$ of
the model's buoyancy term, $\pi_{*,k}-\pi_{*,k-1}=\Delta z\,\Iz\left[g(\rho_Rv_*-1)\right]$. The fallback agrees with the
duals to $O(\Delta z)$. Either way $\pi_*$ is defined up to a constant, which does not affect $A$ (\S\ref{sec:outputs}).

\subsection{FC2D: slots of fixed mass, unknown top pressure}\label{sec:fc-slots}
In a compressible reference column the layer thicknesses are not known in advance. What is known is the pressure as a
function of mass: if $m$ is the mass above a point (per box, from the top), hydrostatic balance gives
$P=\Ptop+gm/L_x$, whatever the arrangement. Only the top pressure $\Ptop$ is unknown, and it is set by the condition that
the parcels fill the box. Take $K$ slots of equal mass $M_k=\Mtot/K$, numbered from the top, with
\begin{equation}
  P_k(\Ptop)=\Ptop+\frac{g}{L_x}\left(k-\tfrac12\right)\frac{\Mtot}{K},\qquad c_{nk}(\Ptop)=h\left(P_k(\Ptop),\vt_n,q_n\right).
  \label{eq:fc-cost}
\end{equation}
$K$ is set by the user. $K=M_z$ gives as many layers as model rows; a finer $K$ improves the vertical resolution of the
reference profile, at a proportional cost.

\paragraph{Energy of a hydrostatic column.} For a column at rest in hydrostatic balance, integrating $\rho g z$ by parts with
$\mathrm{d}P=-\rho g\,\mathrm{d}z$ gives
\begin{equation}
  \int\rho\,(e+\Phi)\,\mathrm{d}A=\int h\,\mathrm{d}m-\Ptop\,L_xL_z-g z_0\Mtot .
  \label{eq:column-energy}
\end{equation}
The discrete form uses the midpoint rule within slots, $\sum_{n,k}x_{nk}c_{nk}$.

\paragraph{The saddle problem.} For a fixed placement, define $G(x,\Ptop)=\sum x_{nk}c_{nk}(\Ptop)-\Ptop L_xL_z$. Then
$\partial G/\partial\Ptop=\sum x_{nk}v(P_k,\vt_n,q_n)-L_xL_z$ is the volume excess, and
$\partial^2G/\partial\Ptop^2=\sum x_{nk}\,\partial_Pv<0$, so $G$ is concave in $\Ptop$ and its maximum is at the $\Ptop$ for
which the placement exactly fills the box. The reference energy is therefore $E_*=E_G-gz_0\Mtot$ with
\begin{equation}
  E_G=\min_x\ \max_{\Ptop}\ G(x,\Ptop)=\max_{\Ptop}\ \min_x\ G(x,\Ptop),
  \label{eq:saddle}
\end{equation}
where the exchange holds because $G$ is linear in the placement $x$, which ranges over a compact convex set, and concave in
$\Ptop$ (Sion's minimax theorem). $G$ is the discrete background energy of the energy notes, restricted to hydrostatic
profiles, and the maximum over $\Ptop$ is their statement that $\int A\,\mathrm{d}m$ is a minimum with respect to $P_*$.
This gives the \textbf{algorithm for FC2D}:
\begin{enumerate}
  \item outer loop: a one-dimensional maximisation over $\Ptop$, by Newton or secant iteration on the volume residual
        $R(\Ptop)=\sum x_{nk}v_{nk}-L_xL_z$, with $R'=-\sum x_{nk}v_{nk}^2/c_{s,nk}^2$ ($c_s$ the sound speed);
  \item inner problem: the rearrangement problem \eqref{eq:tp} with the costs \eqref{eq:fc-cost} at the current $\Ptop$.
\end{enumerate}
Stop on $|R|/(L_xL_z)$ below a tolerance (e.g.\ $10^{-13}$), not after a fixed number of iterations. If the placement is
independent of $\Ptop$ (the separable case of \S\ref{sec:why}, solved by a sort) the inner problem is solved once. Otherwise,
warm-start each inner solve from the previous placement: near convergence it rarely changes.
This is the discrete analogue of the free-surface optimal transport problem of Egan et al.\ (2025), where the surface
pressure plays the role of $\Ptop$.

\paragraph{Geometry of the solution.} Given the converged placement and $\Ptop$, slot $k$ has volume
$V_k=\sum_nx_{nk}v_{nk}$, thickness $\Delta z_k=V_k/L_x$, and heights stacked from the bottom. The edge pressures are
$\Ptop+gk\Mtot/(KL_x)$, and the slot density $\rho_{*,k}=M_k/V_k$ is exactly $-\Delta P_*/(g\,\Delta z_k)$ across the slot,
the discrete form of \eqref{eq:vstar}. Placement heights $\zh_n$ and the sub-slot refinement are as for AN2D.

\paragraph{Initial guess for $\Ptop$.} The mean of the top row of the actual pressure, or the value from the previous snapshot.

\subsection{Dual potentials: the reference pressure}\label{sec:dual}
The dual of \eqref{eq:tp} has a potential $u_n$ per parcel and $\psi_k$ per slot, with
\begin{equation}
  c_{nk}-u_n-\psi_k\ge0\ \ \forall n,k,\qquad\text{with equality wherever }x_{nk}>0 .
  \label{eq:dual}
\end{equation}
So every slot that holds parcel $n$ minimises $c_{nk}-\psi_k$ over $k$, and $u_n$ is that minimum. In \S\ref{sec:theory},
$\zs$ minimises $\Sigma(\vt,q,z,P_*(z,t))$. Comparing the two:
\begin{itemize}
  \item \textbf{AN2D}: $\Sigma(\vt_n,q_n,z_k,P_{*,k})=c_{nk}+\pi_{*,k}$ by \eqref{eq:Sigma}, so
        \begin{equation}
          \pi_{*,k}=-\psi_k+\text{const},\qquad P_{*,k}=P_{R,k}+\rho_{R,k}\,\pi_{*,k}.
        \end{equation}
  \item \textbf{FC2D}: $\Sigma(\vt_n,q_n,z_k,P_{*,k})=c_{nk}+\Phi(z_k)$, so $\psi_k=-\Phi(z_k)+\text{const}$. The duals give
        the heights of the slot centres, which agree with the geometry of \S\ref{sec:fc-slots} to $O(\Delta z)$. The
        pressures are $P_{*,k}=P_k(\Ptop)$.
\end{itemize}
In both models $u_n=\Sigma(\vt_n,q_n,\zs,P_*(\zs,t))$, up to the same constant: the reference static energy in \eqref{eq:A}.
Dual feasibility with complementary slackness is the discrete LNB condition together with the maximal APE condition, in global
form. And $\psi$ is the multiplier of the slot constraints $\sum_nx_{nk}=M_k$, the discrete volume constraint, just as $P_*$ is
the Lagrange multiplier of the volume constraint in the energy notes. For a parcel split between slots $k$ and $k'$, equality
for both slots gives $\psi_k-\psi_{k'}=c_{nk}-c_{nk'}$. In AN2D this is $\pi_{*,k}-\pi_{*,k'}=c_{nk'}-c_{nk}$, the discrete form
of $\partial_z\pi_*=g(\rho_Rv_*-1)$ with $v_*$ the specific volume of the parcel that straddles the interface. Where no parcel
is split across an interface (a degenerate placement), the step of $\psi$ there is only bracketed by the cost differences of the
parcels on either side. This is the discrete counterpart of a jump in $v_*$.

\paragraph{Weak duality and the AE.} The dual objective
\begin{equation}
  D(\psi)=\sum_kM_k\psi_k+\sum_nm_n\min_k\left(c_{nk}-\psi_k\right)
  \label{eq:dual-objective}
\end{equation}
is the background energy of the energy notes for the profile $\psi$, with every parcel at the $\zs$ that maximises its $A$.
In AN2D this holds up to the constant $-\sum_kM_xP_{R,k}\Delta x\Delta z$; in FC2D, $D-\Ptop L_xL_z-gz_0\Mtot$ at the $\Ptop$ of
the costs. Weak duality gives $D(\psi)\le\min_x\sum x_{nk}c_{nk}$ for \emph{every} $\psi$ (and every $\Ptop$). So for any
reference profile, Lorenz or not, $\int A\,\mathrm{d}m$ with the global minimiser $\zs$ is at least the Lorenz APE $E_P-E_*$,
with equality for the optimal duals (cf.\ Tailleux 2013 for a Boussinesq fluid). In FC2D this holds up to the discretisation
mismatch of \S\ref{sec:decisions}. In AN2D with one slot per row it is exact: with $\zs$ restricted to the row centres,
\begin{equation}
  \sum_nm_nA_n=E_P-\left(D(\psi)-\sum_kM_xP_{R,k}\Delta x\Delta z\right),
\end{equation}
because the actual state also has mass $M_k$ in row $k$. With the optimal duals this is $E_P-E_*$ to round-off. Refining $\zs$
between the rows (\S\ref{sec:lnb}) can only increase it.

% =====================================================================================
\section{Solver 1: sorting}\label{sec:sort}
\textbf{Exact iff \eqref{eq:separable} holds} (here: $\gamma=0$). Cost $O(N\log N)$.
\begin{enumerate}
  \item Sort key: the specific volume at a fixed key pressure, $\nu_n=v(P_{\rm key},\vt_n,q_n)$. When $\gamma=0$ any
        $P_{\rm key}$ gives the same order. When $\gamma\neq0$, choose the mid-depth pressure ($P_R$ at $z=L_z/2$ for AN2D,
        the box mean of $P$ for FC2D), since it minimises the largest pressure offset.
  \item Order the parcels by decreasing $\nu$ (lightest first), i.e.\ the top of the column first.
  \item Fill the slots from the top in that order, splitting a parcel across a slot boundary when needed. This gives $x_{nk}$.
  \item FC2D: the placement does not change with $\Ptop$ when the order is exact; run the outer Newton loop of
        \S\ref{sec:fc-slots} once on this placement. With $\gamma\neq0$, the order can in principle be refreshed at the final
        pressures, which is not a minimisation; prefer solvers 2 or 3 if the sort is found to be inadequate.
\end{enumerate}
At parcel resolution ($K=N$, a slot per parcel in sort order) this is the construction of Winters et al.\ (1995) and of
the earlier \code{FCSCEnergetics.jl}.

\paragraph{Duals and certificate.} The positive entries of the sorted placement form a staircase. Solve $u_n+\psi_k=c_{nk}$ along
it from the top, in $O(N+K)$; where no parcel is split across a slot boundary, complete the staircase with a zero entry. Then
check \eqref{eq:dual} for all $NK$ parcel--slot pairs (an $O(NK)$ pass). If it holds, the sorted placement is the exact optimum, not only
2-optimal, and $\psi$ is the reference profile of \S\ref{sec:dual}. If it fails, the violating pairs show where to repair
(\S\ref{sec:repair}) or which pairs to add to the banded linear programme (\S\ref{sec:exact}).

\paragraph{Pros.} Fast and simple, with no dependency. It is the baseline against which every other solver is tested:
with $\gamma=0$ all solvers must agree with it to round-off.
\paragraph{Cons.} Approximate when $\gamma\neq0$. The certificate, or the duality gap of \S\ref{sec:repair}, measures the
error.

% =====================================================================================
\section{Solver 2: greedy filling with repair}\label{sec:greedy}
Approximate for any equation of state, with no dependency. Each method below evaluates the specific volume of the candidates
at the pressure of the level being filled, so it accounts for thermobaricity.

\subsection{Top-down filling}
For $k=1,\dots,K$ from the top: evaluate $v(P_k,\vt_n,q_n)$ for the parcels not yet fully placed, and fill slot $k$ with the
largest values (lightest first), splitting the last parcel. Cost: $O(NK)$ evaluations of $v$ (about $4\times10^6$ at
$256\times128$, $K=128$). This is the ``greedy'' algorithm of Huang (2005) and Wong et al.\ (2016), and the discrete
counterpart of the top-down construction of Saenz et al.\ (2015). \emph{Bottom-up} is the same from $k=K$ upwards,
taking the smallest $v$ first. Harris \& Tailleux (2018) found top-down to be the least reliable heuristic for moist air,
and bottom-up to systematically miss part of the APE there.

\subsection{Divide and conquer}
Stansifer et al.\ (2017). On a slot range $[k_1,k_2]$ with its set of parcel mass: split it at the slot boundary
$k_m$ that halves the mass, evaluate $v$ at the boundary pressure, and send the lightest parcels (by mass, splitting one
parcel at the cut) to $[k_1,k_m]$ and the rest to $[k_m+1,k_2]$. Recurse until single slots. Cost: $O(N\log K)$ evaluations.
This was the best speed/accuracy compromise in Harris \& Tailleux (2018), but it still had $>50\%$ error in 11\% of their
cases and occasionally gave negative APE, so the repair below is not optional.

\subsection{Repair (pairwise exchanges, 2-opt) and the energy check}\label{sec:repair}
An exchange of mass $\delta>0$ between parcel $n$ in slot $k$ and parcel $l$ in slot $k'$ changes the cost by
\begin{equation}
  \delta\,\left(c_{nk'}+c_{lk}-c_{nk}-c_{lk'}\right).
\end{equation}
Whenever the bracket is negative, exchange $\delta=\min(x_{nk},x_{lk'})$. Repeat over all pairs of slots with
$|k-k'|\le B$ (a band; thermobaric misorderings are local, so a small $B$ such as $2$--$4$ should suffice; check by
increasing $B$) until no improving exchange is found. The result is a local minimum, not necessarily the global one
(Hieronymus \& Nycander 2015 give examples where 2-opt fails). For FC2D the repair runs inside the outer $\Ptop$ loop.

\textbf{Energy check (mandatory for every approximate solver):} the reference energy must not exceed the energy of the
actual state, $E_*\le E_P$. For AN2D with one slot per row the actual state is a feasible placement with exactly the same discrete
energy, so a violation is a solver failure and must raise an error. For FC2D (and AN2D with $s>1$) the two energies are
computed on different grids (\S\ref{sec:decisions}), so allow a tolerance set by the discretisation.

\textbf{Duality gap:} any feasible placement $x$ gives $E[x]\ge E_*$, and any $\psi$ gives $D(\psi)\le E_*$ (up to the constants of
\S\ref{sec:dual}). With $\psi$ from the placement (the staircase of \S\ref{sec:sort}) or from the slot means
(\S\ref{sec:an-slots}), $E[x]-D(\psi)$ bounds the error of the approximate solver at the cost of one $O(NK)$ pass.

\paragraph{Pros.} No dependency; handles thermobaricity; near-exact in practice after repair.
\paragraph{Cons.} No optimality guarantee; heuristics without repair can give negative APE; slower than the sort.

% =====================================================================================
\section{Solver 3: exact rearrangement}\label{sec:exact}
Solves \eqref{eq:tp} exactly for any equation of state. Its outputs are the placement $x$ and the dual potentials
$(u_n,\psi_k)$, which certify optimality by \eqref{eq:dual} and give the reference profile (\S\ref{sec:dual}).
\subsection{Banded linear programme with a dual certificate (recommended)}
The full problem has $NK\approx4\times10^6$ parcel--slot pairs at $256\times128$. Because parcels move only a few slots away from their
sorted position, solve it on a band and certify the result:
\begin{enumerate}
  \item Compute the sort (\S\ref{sec:sort}); let $k_n$ be the slot of parcel $n$ in the sorted placement.
  \item Solve the restricted problem with the pairs $|k-k_n|\le B$ only (start with $B=2$), as a linear programme or a min-cost
        flow, and read the duals $(u_n,\psi_k)$ from the solver.
  \item Check \eqref{eq:dual} for \emph{all} $NK$ pairs (an $O(NK)$ pass over $c_{nk}$, cheap). If every reduced cost is
        nonnegative, the placement is the exact optimum of the full problem. Otherwise add the violating pairs (or double $B$) and
        resolve.
\end{enumerate}
\textbf{Solver}: an external LP solver, HiGHS (e.g.\ through JuMP), which returns the placement and the duals of the equality
constraints; a min-cost-flow / network simplex package is an alternative. For equal-mass parcels and $K=N$ the problem is the
linear assignment problem (Hungarian or Jonker--Volgenant), but that is $O(N^3)$ and not needed here. Su \& Ingersoll (2016) report min-cost flow $10^3$--$10^5$ times faster than linear
assignment on the same problem, by exploiting $K\ll N$, which is the structure used here.

\subsection{Alternative: dual ascent}
With an LP solver this is not needed, but it is a dependency-free route and it makes the link with \S\ref{sec:theory}
explicit. The dual of \eqref{eq:tp}, after eliminating $u$, is the maximisation of the concave piecewise-linear $D(\psi)$
\eqref{eq:dual-objective} in only $K$ variables, with supergradient $M_k-(\text{mass of the parcels whose minimiser is }k)$.
It is the discrete form of the variational principle of \S\ref{sec:theory}: for a given profile every parcel takes the level
that maximises its $A$, and the profile is adjusted until the volume constraint holds; the supergradient is the volume residual
of each slot. It is also the discrete counterpart of the semi-discrete optimal transport (Laguerre-cell) formulation of Egan
et al.\ (2025). It can be maximised by a bundle or subgradient method, or by entropic smoothing (log-domain Sinkhorn with
decreasing regularisation). The placement is read off from the minimisers. Parcels with ties must be split by a small flow
among the tied slots. This needs more numerical care than the LP route.

\paragraph{Pros.} The true Lorenz state, with a certificate and the reference pressure; handles unequal parcel masses and any
equation of state; the band keeps it close to the cost of the sort.
\paragraph{Cons.} An external dependency (HiGHS). It is slower than the sort, and for FC2D it is repeated for each $\Ptop$
iteration (warm-started).

% =====================================================================================
\section{Solver 4: histogram (PDF) reduction}\label{sec:histogram}
Tseng \& Ferziger (2001), generalised to two tracers by Saenz et al.\ (2015). This is a \emph{reduction} rather than a solver:
it replaces the $N$ parcels by $N_b$ bins before applying one of solvers 1--3.
\begin{enumerate}
  \item Bin the parcels in $(\vt,q)$ (or in the sort key $\nu$ if separable). For each nonempty bin $b$ store its mass
        $m_b=\sum_{n\in b}m_n$ and representative properties: the mass-weighted means $\bar\vt_b,\bar q_b$.
  \item Treat the bins as parcels and solve \eqref{eq:tp} with solver 1, 2 or 3.
  \item Give the parcels their reference height by the LNB equation against the resulting reference profile
        (\S\ref{sec:lnb}), which gives a continuous $\zs$ for each parcel rather than one per bin.
\end{enumerate}
In the separable case the profile follows directly from cumulative mass: a bin at sort key $\nu$ sits at the height where the
mass of lighter fluid above it equals the mass of all bins with larger $\nu$ (Tseng \& Ferziger 2001, eq.\ 6, with mass in
place of volume).
\paragraph{Pros.} $O(N+N_b)$; a smooth profile with a resolution set by the bins rather than the grid; the bin counts are
reusable across diagnostics.
\paragraph{Cons.} Averaging $(\vt,q)$ inside a bin with a nonlinear equation of state biases the energy (cabbeling), so bins
must be fine. A resting, stably stratified state no longer gives zero APE exactly. At $N\approx3\times10^4$ the saving over
the sort is small; the benefit is for much larger grids or long time series.

% =====================================================================================
\section{Reference heights: the LNB equation}\label{sec:lnb-sec}
Given a reference profile, Lorenz (\S\ref{sec:lorenz}) or horizontal mean (\S\ref{sec:mean}), each parcel gets its $\zs$ by
minimising $\Sigma(\vt,q,z,P_*(z,t))$, as in \S\ref{sec:theory}.

\subsection{Profile representation}\label{sec:profile}
A profile is known at $K$ levels (slot edges or centres, or model rows) and must be interpolated to solve the LNB equation.
Interpolate the \emph{pressures} and derive the specific volumes from the interpolant. Then $A$, the LNB residual, $v_*$ and
$N_*^2$ all come from one profile, and the $\zs$ found below is an exact minimiser of the $\Sigma$ that enters $A$:
\begin{itemize}
  \item \textbf{AN2D}: interpolate $P_R$ and $\pi_*$ between the row centres with $C^1$ cubic Hermite interpolants (for
        $P_R$ with the node slopes $-g\rho_{R,k}$, so that it matches the model's $P_R$ and $\rho_R$ at the centres). Set
        $\rho_R(z)=-P_R'(z)/g$ and
        $v_*(z)=\left(\pi_*'(z)/g+1\right)/\rho_R(z)$ by \eqref{eq:vstar};
  \item \textbf{FC2D}: interpolate $P_*$ in height through the slot-edge pressures (Lorenz) or the row values (horizontal
        mean) with a monotone $C^1$ cubic, and set $\rho_*(z)=-P_*'(z)/g$, $v_*=1/\rho_*$. Do not interpolate $\rho_{*,k}$
        independently.
\end{itemize}
Beyond the first and last levels, extend the pressures linearly (constant $v_*$). A piecewise-linear interpolant also works
with the global minimum of \S\ref{sec:lnb}, but it gives a piecewise-constant $v_*$, and the minimum may then sit at a level,
where the residual changes sign without vanishing.

\subsection{The LNB equation and the choice of root}\label{sec:lnb}
For a parcel $(\vt,q)$, write $S(z)=\Sigma(\vt,q,z,P_*(z,t))$ for its static energy along the profile, \eqref{eq:Sigma}.
Its reference height is the global minimiser
\begin{equation}
  \zs=\arg\min_{0\le z\le L_z}S(z)
  \label{eq:zstar}
\end{equation}
(\S\ref{sec:theory}). At an interior minimum $r(\zs)=0$, with the \emph{LNB residual}
\begin{equation}
  r(z)=-\frac{\mathrm{d}S}{\mathrm{d}z}=\begin{cases}
    g\,\rho_R(z)\left(v(P_R(z),\vt,q)-v_*(z)\right) & \text{AN2D},\\[2pt]
    g\left(\rho_*(z)\,v(P_*(z),\vt,q)-1\right) & \text{FC2D},
  \end{cases}
\end{equation}
and there $N_*^2=-\mathrm{d}r/\mathrm{d}z\ge0$. $r>0$ means the parcel is lighter than the reference at $z$. $N_*^2$ depends
on the parcel through $(\vt,q)$ and not only on height (Tailleux 2013, Appendix A).
\begin{enumerate}
  \item \textbf{Coarse search:} evaluate $S$ at the $K$ levels of the profile and at the walls (cheap: $K\approx128$
        evaluations per parcel) and list its local minima. At the levels, $S=c_{nk}+\pi_{*,k}$ (AN2D) or
        $h(P_{*,k},\vt,q)+\Phi(z_k)$ (FC2D). For a Lorenz profile from the duals this is $c_{nk}-\psi_k$ up to a constant
        (\S\ref{sec:dual}), the same pass as the certificate of \S\ref{sec:exact}.
  \item \textbf{Refine} each local minimum whose value is within a tolerance of the smallest. On the intervals next to it,
        find the root of $r$ with $\mathrm{d}r/\mathrm{d}z<0$ by Brent's method on the interpolated profile, to a tolerance
        well below $\Delta z$. Keep the smallest refined $S$. A wall is the minimum if $r$ has the outward sign there ($r\ge0$
        at the top, $r\le0$ at the bottom).
  \item \textbf{Record} $N_*^2$; the parcels at a wall, where $r(\zs)\ne0$ (evaluate their budget terms with the general
        forms of the energy notes, $\Sigma_P(\vt,q,\zs,P_*(\zs,t))$ rather than $v_*(\zs,t)$); and the parcels with two minima
        within the tolerance, where $\zs$ can jump between them.
\end{enumerate}

\paragraph{Why the global minimum.} It maximises the parcel's $A$ at fixed $P_*$ (\S\ref{sec:theory}), so it depends on
$(\vt,q,t)$ only, not on the parcel's actual height, and $\Pi_2=S(z)-S(\zs)\ge0$. $S(\zs)$ is continuous in $(\vt,q,t)$ even
where $\zs$ jumps between two minima. ($\Sigma_\vt$ and $\Sigma_q$ at $\zs$, which enter the heating terms of the energy notes,
do jump there.) For the Lorenz profile it reproduces the levels of the placement, and for any profile it gives the
bound of \S\ref{sec:dual}. Tailleux (2013, \S3.1) instead takes the root nearest to the parcel's height $z$: above $z$ if
$r(z)>0$, below if $r(z)<0$, or the wall if there is none. This is also positive definite, but then the reference height
depends on $z$, and $A$ jumps when a parcel crosses a local maximum of $S$, so it is not used here. The two rules agree when
$S$ has a single local minimum.

Multiple minima arise for the Lorenz profile only through the pressure dependence of the parcel ordering ($\gamma\ne0$; the
overlap region of Saenz et al.\ 2015). For the horizontal mean they also arise where the mean profile is statically unstable,
e.g.\ during overturning in Kelvin--Helmholtz billows. Count and record parcels with multiple minima and with a wall minimum.

\subsection{Against the Lorenz profile}
For the Lorenz profile, $\zs$ agrees with the placement height $\zh$ to the slot resolution. In slots that hold a
mixture of parcels $v_*$ is an average, so the two can differ by up to a slot thickness there. The budgets use $\zs$: it
satisfies the LNB condition for the profile used in $A$, so the $\partial_t\zs$ terms vanish (\S\ref{sec:theory}). The
placement height $\zh$ satisfies the volume constraint exactly but the LNB condition only to $O(\Delta z)$; for $\zs$ it is the other
way round, and the LNB condition is the one the local budgets need. The LNB $\zs$ is also a single height for a parcel split
between slots, so the fractions $x_{nk}$ are not needed for $A$.

% =====================================================================================
\section{Horizontal-mean reference state}\label{sec:mean}
This reference state is \emph{not} a rearrangement of the fluid: its $P_*$ does not satisfy the volume constraint
\eqref{eq:volume}. The local AE and its local budgets in the energy notes hold for any $P_*(z,t)$ with the $\zs$ of
\S\ref{sec:lnb}. But the reference pressure tendency term $R_A$ no longer integrates to zero, and $\int A\,\mathrm{d}m$ exceeds
the Lorenz APE (\S\ref{sec:dual}; Tailleux 2013). It is cheap and needs no sort.

\subsection{The profiles}\label{sec:mean-profiles}
The profile is specified by $P_*(z,t)$ alone, hydrostatic with the horizontal mean of the quantity in the model's hydrostatic
balance. In FC2D this is the density, the usual choice in the literature (Tailleux 2013; Tailleux 2018, \S4.2):
\begin{itemize}
  \item \textbf{AN2D.} The thermodynamics stays at $P_R$. $\pi_*$ is hydrostatic with the row mean
        $v_{*,k}=\frac1{M_x}\sum_iv(P_{R,k},\vt_{ik},q_{ik})$, by the rule of \S\ref{sec:an-slots} with $s=1$. This is the row
        mean of the buoyancy $b=g(\rho_Rv-1)$, which is linear in $v$, as the FC2D balance is linear in $\rho$.
  \item \textbf{FC2D.} $P_*$ is hydrostatic with the row-mean density $\bar\rho_k=\frac1{M_x}\sum_i m^\rho_{ik}$ by the
        model's discrete rule $P_{*,k}-P_{*,k-1}=-g\Delta z\,\Iz\bar\rho$, anchored to the row-mean pressure at one level.
        The cells have equal volume, so this keeps the mass of each row, and $\rho_*=\bar\rho$.
\end{itemize}
In both, $v_*$ is still derived from the interpolant of $P_*$ (or $\pi_*$) of \S\ref{sec:profile}. Hermite node slopes from the
row means ($-g\bar\rho_k$, or $g(\rho_{R,k}v_{*,k}-1)$) make it equal to the row mean at the centres. The row mean of the
pressure itself is not used. It contains the non-hydrostatic part of the mean pressure, so its implied density
$-\partial_zP_*/g$ is not $\bar\rho$ and it does not keep the row masses.

% =====================================================================================
\section{Time derivative of the reference pressure}\label{sec:dtpstar}
The reference pressure tendency terms $R_{\Pi_1}$, $R_{\Pi_2}$ and $R_A$ of the energy notes need $\partial_tP_*(z,t)$ on the
model rows and at each parcel's $\zs$, from the same interpolant as $P_*$ (\S\ref{sec:profile}). By default it is computed
from the diabatic tendencies of a single snapshot (\S\ref{sec:dtpstar-b}); differences of successive snapshots
(\S\ref{sec:dtpstar-a}) serve as a check. Both hold for any equation of state: they use only $v(P,\vt,q)$ and its derivatives,
and keep the compressibility $v_P$ wherever it enters.

The global identity $\int R_A\,\mathrm{d}V=0$ of the energy notes uses only the volume constraint \eqref{eq:volume} at time
$t$, with $f=\partial_tP_*$. It therefore holds however $\partial_tP_*$ is estimated, to the accuracy with which $\zs$ satisfies
the volume constraint (\S\ref{sec:lnb-sec}). The estimate affects only the local distribution of $R_A$ and, in FC2D, the
exchange $\int R_{\Pi_1}\,\mathrm{d}V$ between $\Pi_1$ and $\Pi_2$.

\subsection{From the diabatic tendencies (default)}\label{sec:dtpstar-b}
The Lorenz reference state depends only on the distribution of mass over $(\vt,q)$. Advection rearranges the fluid without
changing that distribution, so $P_*$ changes only through the diabatic rates of the fluid elements (Winters et al.\ 1995;
Winters \& D'Asaro 1996),
\begin{equation}
  \dot\vt=\frac{D\vt}{Dt}=\frac{\sigma_\vt-\nabla\cdot\mathbf J_\vt}{\rho},\qquad
  \dot q=\frac{Dq}{Dt}=-\frac{\nabla\cdot\mathbf J_q}{\rho},
\end{equation}
which the model provides (\code{conservative\_variable\_production}, \code{div\_consflux}, \code{div\_qflux}; $\rho=\rho_R$ in
AN2D). For a parcel at pressure $P$, let
\begin{equation}
  \dot v=v_\vt\,\dot\vt+v_q\,\dot q
\end{equation}
be the rate of change of its specific volume at fixed pressure, with $v_\vt=\partial_Ph_\vt$ and $v_q=\partial_Ph_q$ (or a
directional derivative of $v(P,\vt,q)$).

\paragraph{FC2D.} In the mass coordinate $m$ of \S\ref{sec:fc-slots} (the mass above a level), $P_*(m,t)=\Ptop(t)+gm/L_x$. Let $v_*(m,t)$ be the
reference specific volume and $m_*(\vt,q,t)$ the reference mass coordinate of a parcel. Follow a fluid element, whose mass is
fixed, and differentiate the LNB condition $v(P_*(m_*,t),\vt,q)=v_*(m_*,t)$ along it. With $\partial_m=-(v_*/L_x)\,\partial_z$
and $N_*^2=g\rho_*^2\left(v_*\partial_zv_*+gv_P\right)$ (energy notes, ``Maximal APE condition''),
\begin{equation}
  \frac{N_*^2}{g\rho_*^2L_x}\,\frac{Dm_*}{Dt}=\partial_tv_*|_m-\dot v-v_P\,\partial_t\Ptop ,
\end{equation}
with $\dot v$, $v_P$ and $N_*^2$ those of the parcel at its reference level. The mass above the level $m$ is $m$ itself at all
times, so the mass-weighted mean of $Dm_*/Dt$ over the parcels at a level vanishes. Hence
\begin{equation}
  \partial_tv_*|_m=\frac{\left\langle\left(\dot v+v_P\,\partial_t\Ptop\right)/N_*^2\right\rangle_m}
                          {\left\langle1/N_*^2\right\rangle_m},
  \label{eq:dtvstar}
\end{equation}
where $\langle\cdot\rangle_m$ is the mass integral over the parcels whose reference level is $m$. The box stays full,
$\int_0^{\Mtot}\partial_tv_*|_m\,\mathrm{d}m=0$, which fixes $\partial_t\Ptop$: a heated column expands, and the lid compresses
it back. With $z(m,t)=L_x^{-1}\int_m^{\Mtot}v_*\,\mathrm{d}m'$, the tendency at fixed height is
\begin{equation}
  \partial_tP_*|_z=\partial_t\Ptop+g\rho_*\,\partial_tz|_m,\qquad
  \partial_tz|_m=\frac1{L_x}\int_m^{\Mtot}\partial_tv_*|_{m'}\,\mathrm{d}m' .
  \label{eq:dtPz}
\end{equation}

\paragraph{AN2D.} The mass per unit height is fixed, so the levels do not move and there is no lid. The same steps with
$v(P_R(\zs),\vt,q)=v_*(\zs,t)$ give $(N_*^2/g\rho_R)\,D\zs/Dt=\dot v-\partial_tv_*$, hence
\begin{equation}
  \partial_tv_*(z,t)=\frac{\left\langle\dot v/N_*^2\right\rangle_z}{\left\langle1/N_*^2\right\rangle_z},\qquad
  \frac{\partial}{\partial z}\left(\partial_t\pi_*\right)=g\rho_R\,\partial_tv_*,\qquad \partial_tP_*=\rho_R\,\partial_t\pi_*,
  \label{eq:dtvstar-an}
\end{equation}
with $\dot v$ and $N_*^2=g\rho_R\left(\partial_zv_*+g\rho_Rv_P\right)$ at $(P_R(\zs),\vt,q)$. The constant in $\partial_t\pi_*$ is
free, since $R_A$ does not depend on it.

\paragraph{The weights.} The reference state follows the $N_*^{-2}$-weighted mean of the parcels at a level. A parcel that is
weakly stratified at its level moves further for a given change of $v$, so it displaces more of its neighbours. The weights
differ between the parcels of a level only through $v_P$, i.e.\ through the pressure dependence of $v$ (thermobaricity), and
they differ most where $N_*^2$ is small. The plain mean is exact only when $v_P$ is the same for all of them, e.g.\ for
\eqref{eq:separable}. It is also what a derivative at fixed placement gives, since a fixed placement does not see the parcels reorder
within a level. So the weights must be kept for a general equation of state (\S\ref{sec:tests}, test~\ref{test:weights}).

\paragraph{Discrete form.} The parcels at a level are those of a slot, with the weights of the placement. For parcel $n$ in slot $k$,
evaluate $\dot v_{nk}$ and $v_{P,nk}$ at the slot pressure ($P_{R,k}$ in AN2D, $P_k$ in FC2D), and
\begin{equation}
  N^2_{nk}=g\rho_k\left(\partial_zv_*+g\rho_k\,v_{P,nk}\right),\qquad w_{nk}=\frac{x_{nk}}{N^2_{nk}},\qquad
  \langle a\rangle_k=\frac{\sum_nw_{nk}\,a_{nk}}{\sum_nw_{nk}},
\end{equation}
with $\rho_k=\rho_{R,k}$ (AN2D) or $\rho_{*,k}$ (FC2D), and $\partial_zv_*$ at $z_k$ from the profile of \S\ref{sec:profile}.
\begin{itemize}
  \item \textbf{AN2D}: $\partial_tv_{*,k}=\langle\dot v\rangle_k$, and $\partial_t\pi_*$ by integrating \eqref{eq:dtvstar-an}
        between the slot centres with the rule used for $\pi_*$ (\S\ref{sec:an-slots}).
  \item \textbf{FC2D}:
        \begin{equation}
          \partial_t\Ptop=-\frac{\sum_kM_k\langle\dot v\rangle_k}{\sum_kM_k\langle v_P\rangle_k},\qquad
          \partial_tv_{*,k}=\langle\dot v\rangle_k+\langle v_P\rangle_k\,\partial_t\Ptop .
        \end{equation}
        The slot edges move at $\partial_tz_e=L_x^{-1}\sum_{k\text{ below }e}M_k\,\partial_tv_{*,k}$, which vanishes at both
        walls, and the edge pressures at $\partial_t\Ptop$. $\partial_tP_*$ at fixed height is the derivative of the
        interpolant of \S\ref{sec:profile} with respect to these node motions; at the edges this is \eqref{eq:dtPz}.
\end{itemize}
If $N^2_{nk}\le0$ for a parcel of a slot (a marginally or discretely unstable level), the weights are undefined there: use
$w_{nk}=x_{nk}$ in that slot and flag it. The tendency is that of the slot-level profile, so it matches the dual profile of
\S\ref{sec:dual} to $O(\Delta z)$.

\paragraph{Horizontal-mean reference state.} Its profile is a row average, not a function of the distribution of $(\vt,q)$,
so it also changes under advection. Differentiate its definition (\S\ref{sec:mean-profiles}) with the model's full
tendencies. In AN2D, $\partial_tv_{*,k}$ is the row mean of $v_\vt\partial_t\vt+v_q\partial_tq$ at $P_{R,k}$. In FC2D,
$\partial_t\bar\rho_k$ is the row mean of the mass tendency, $\partial_tP_*$ is hydrostatic with it by the same discrete rule,
and the anchor moves with the row mean of $\partial_tP=(\partial_tv-v_\vt\partial_t\vt-v_q\partial_tq)/v_P$ at its level.

\subsection{From differences of successive snapshots (check)}\label{sec:dtpstar-a}
Compute the reference state at each snapshot, warm-starting the solver from the previous placement. Evaluate $P_*$ on the model
rows, take centred differences in time, and evaluate the result at $\zs$ with the interpolant of the middle snapshot.
\begin{itemize}
  \item In AN2D the constants in $\pi_*$ need not agree between snapshots, since $R_A$ is unchanged under
        $\pi_*\to\pi_*+c(t)$.
  \item Where the placement is degenerate the duals are not unique (\S\ref{sec:dual}), and a solver may return a different value at
        each snapshot. Difference the slot-mean profile there, or make the choice of dual reproducible.
  \item This is the actual time derivative of the diagnosed reference state, including the effect of numerical mixing. A
        budget of $A$ differenced between the same snapshots closes against it, to $O(\Delta t^2)$.
  \item The cell-based reference state is not exactly invariant under discrete advection: as features cross cells, the
        histogram of cell values changes slightly, at about $|\mathbf u|/\Delta x$. Use a snapshot interval
        $\Delta t\gg\Delta x/|\mathbf u|$, or average, so that this does not alias into $\partial_tP_*$. The interval must
        still be short compared with the diabatic evolution of $P_*$.
\end{itemize}

\paragraph{Using both.} The estimate from the diabatic tendencies excludes, by construction, any change of the reference state
caused by the discretisation of advection; the snapshot differences include it. Their difference therefore measures the
spurious (numerical) change of the reference state.
This is the local counterpart of using the background potential energy to measure spurious mixing (Il{\i}cak et al.\ 2012),
and is useful for comparing the schemes of the models.

% =====================================================================================
\section{Outputs}\label{sec:outputs}
Each solver returns, for one snapshot, what the energy notes need (\S\ref{sec:theory}):
\begin{itemize}
  \item the \textbf{reference profile}, in the representation of \S\ref{sec:profile}. AN2D: $\pi_{*,k}$ on the slot centres (so
        $P_*=P_R+\rho_R\pi_*$). FC2D: the slot edges and centres $z_k$, $P_{*,k}$, $\rho_{*,k}=M_k/V_k$ and $\Ptop$. Both also
        return $P_*$ and $v_*$ at the model rows $z_j$ and faces, from the same interpolant;
  \item per parcel: $\zs$ (the global minimiser of \S\ref{sec:lnb}); the reference static energy
        $\Sigma(\vt,q,\zs,P_*(\zs,t))$; the pressure at which the thermodynamics is evaluated there ($P_*(\zs,t)$ in FC2D,
        $P_R(\zs)$ in AN2D, with $\pi_*(\zs,t)$); $N_*^2$; the placement height $\zh$; and flags for split, wall and
        multiple-minimum parcels;
  \item the \textbf{reference pressure tendency} $\partial_tP_*$ on the model rows and at each $\zs$, from the diabatic
        tendencies (\S\ref{sec:dtpstar-b}), with $\partial_tv_{*,k}$, $\partial_t\Ptop$ (FC2D) and the slots where the weights
        were replaced. Given a time series, also from differences of successive snapshots (\S\ref{sec:dtpstar-a}), as a
        check;
  \item the \textbf{potential energies} $E_P$ of the actual state and $E_*$ of the reference state, in the same
        discretisation (AN2D: $E[x]$; FC2D: \eqref{eq:column-energy}), and $\int A\,\mathrm{d}m$ from \eqref{eq:A} as a check
        (\S\ref{sec:dual});
  \item for the exact rearrangement, the duals and the certificate status; for the iterative ones, the iteration counts and
        final residuals.
\end{itemize}
In AN2D the constant in $\pi_*$ is free at each time: $A$, $\zs$ and the reference pressure tendency term $R_A$ of the energy
notes are unchanged under $\pi_*\to\pi_*+c(t)$. In FC2D, $P_*$ is absolute (it is fixed by $\Ptop$). FC2D energies and APE carry the $O(\Delta z^2)$ grid mismatch accepted in
\S\ref{sec:decisions}.

% =====================================================================================
\section{Summary of the methods}
\begin{center}\small
\begin{tabular}{p{2.8cm}p{3.2cm}p{2.3cm}p{5.8cm}}
\hline
Method & Exact? & Cost & Use \\ \hline
Sort & iff $\gamma=0$ & $O(N\log N)$ & baseline, default when certified (\S\ref{sec:sort}) \\
Greedy + repair & no (local minimum) & $O(NK)$ & fallback if the exact solver is too slow \\
Exact rearrangement & yes, certified & band LP & ground truth; validation of 1, 2, 4; $P_*$ from the duals \\
Histogram & to bin resolution & $O(N+N_b)$ & large grids, long time series \\
Mean + LNB & not a rearrangement & $O(NM_z)$ & cheap, local reference; APE above Lorenz \\ \hline
\end{tabular}
\end{center}
Order of implementation: the sort with its certificate and the LNB (\S\ref{sec:lnb-sec}) first, with the mean profile (simple
and cheap); then the exact rearrangement, as the reference solver and to measure the error of the sort with
$\gamma\neq0$; then greedy+repair and the histogram only if the exact solver is too slow for the intended use.

% =====================================================================================
\section{Decisions}\label{sec:decisions}
\paragraph{Settled with the user.}
\begin{itemize}
  \item \textbf{Horizontal-mean profile}: hydrostatic with the row mean of the density (FC2D) or of the specific volume at
        $P_R$ (AN2D), \S\ref{sec:mean-profiles}.
  \item \textbf{Exact solver}: an external LP solver (HiGHS) is acceptable in \code{TraceableEnergetics}, \S\ref{sec:exact}.
  \item \textbf{Slot resolution} $K$: an input set by the user (\S\ref{sec:an-slots}, \S\ref{sec:fc-slots}).
  \item \textbf{Precision}: any floating-point type $F$, chosen by the user (Float128 benchmarks are planned), so no step of
        the computation may be fixed to Float64.
  \item \textbf{Lorenz profile}: from the duals when the solver provides them, otherwise from slot means
        (\S\ref{sec:an-slots}).
  \item \textbf{Reference height for the budgets}: the LNB $\zs$, not the placement height $\zh$ (\S\ref{sec:lnb-sec}).
  \item \textbf{Choice of root}: the global minimiser (\S\ref{sec:lnb}), not the nearest root of Tailleux (2013).
  \item \textbf{Reference pressure tendency}: from the diabatic tendencies ($N_*^{-2}$ weights, any equation of state) by
        default, with differences of successive snapshots as a check (\S\ref{sec:dtpstar}).
\end{itemize}

\paragraph{Accepted limitation: the FC2D reference grid.} In AN2D the reference column is built on the model's own rows (one slot per row),
and the actual state is one of the candidate rearrangements, scored with exactly the model's energy. A fluid that is already at
rest and stably stratified is then its own reference state, with zero APE to round-off. FC2D cannot be set up this way. In a
compressible column, the mass that fits in a row of fixed height depends on which parcels are put there, so the capacities of
the rows are not known until the reference state is. The reference column is therefore built from layers of fixed mass, whose
thicknesses come out of the calculation (\S\ref{sec:fc-slots}). These layers do not line up with the model rows. Their energy
is also computed by a different formula, \eqref{eq:column-energy}, rather than the model's sum over cells with its discrete
hydrostatic balance. Both are second-order accurate, but they do not give the same number: a fluid already in its Lorenz state
gets an APE of order $\Delta z^2$, of either sign. This decreases as the model grid is refined, but not by refining $K$ alone,
because part of it is the model's own discretisation error. Removing it would need slots on the model rows with unknown
masses, a free-surface rearrangement problem that is harder to solve. \textbf{The $O(\Delta z^2)$ mismatch is accepted} (decision
of the user). It is present in every FC2D result: the energy check of \S\ref{sec:repair} uses a matching tolerance, test~1 of
\S\ref{sec:tests} measures it, and FC2D diagnostics should report it alongside the APE.

% =====================================================================================
\section{Tests}\label{sec:tests}
\begin{enumerate}
  \item \textbf{Resting stable state}, AN2D with one slot per row: a horizontally uniform, stably stratified state gives
        $x$ = identity, $\zh=z$ and $E_*=E_P$ to round-off, for every solver. With the dual profile and $\zs$ restricted to
        the row centres, $A=0$ to round-off for every parcel. For FC2D (and AN2D with $s>1$) the same holds up to
        $O(\Delta z^2)$, converging as the model grid is refined.
  \item \textbf{Separable equation of state} ($\gamma=0$): sort, greedy+repair and exact rearrangement give the same $E_*$ to
        round-off on random and Kelvin--Helmholtz snapshots, and the sort certificate holds.
  \item \textbf{Thermobaric} ($\gamma\neq0$): $E_{*,\rm exact}\le E_{*,\rm greedy+repair}\le E_{*,\rm sort}$ (the middle
        inequality holds when the repair starts from the sort). Report the relative gaps.
  \item \textbf{Bounds}: $E_*\le E_P$ (AN2D with $s=1$ to round-off, otherwise within the discretisation tolerance). For any
        profile (Lorenz from any solver, horizontal mean) with $\zs$ the global minimiser, $\int A\,\mathrm{d}m\ge
        E_P-E_{*,\rm exact}$, with equality to round-off in AN2D ($s=1$) for the optimal duals and $\zs$ on the row centres.
  \item \textbf{Conservation}: $\sum_kx_{nk}=m_n$, $\sum_nx_{nk}=M_k$; FC2D $|R(\Ptop)|$ below tolerance and
        $\sum_k\Delta z_k=L_z$.
  \item \textbf{Permutation invariance}: shuffling the parcels (e.g.\ shifting the field periodically in $x$) leaves $E_*$
        and the profile unchanged.
  \item \textbf{Gauge}: changing $z_0$ shifts $E_*$ and $E_P$ by the same constant; the reference heights and $A$ are
        unchanged. In AN2D, $\pi_*\to\pi_*+c$ leaves $\zs$ and $A$ unchanged.
  \item \textbf{LNB}: $r(\zs)=0$ to tolerance at interior minima, with $N_*^2\ge0$; $A\ge0$ for every parcel ($\Pi_1\ge0$
        by the concavity of $\Sigma$ in $P$, $\Pi_2\ge0$ by construction). For the Lorenz profile, $\zs=\zh$ up to the slot
        resolution in unmixed slots.
  \item \textbf{Certificate}: the duals of the exact solver satisfy \eqref{eq:dual} for all parcel--slot pairs.
  \item\label{test:weights} \textbf{Reference pressure tendency} from the diabatic tendencies:
        \begin{enumerate}
          \item with $v_P$ the same for all parcels of a level ($\gamma=0$), it equals the finite-difference derivative of the
                reference state for a small diabatic step;
          \item with strong thermobaricity, e.g.\ two density-compensated families of parcels with different $v_P$, heated
                differently: the finite-difference derivative of the exact solver, which includes the reordering, matches
                the weighted mean, and not the plain one;
          \item FC2D: $\sum_kM_k\,\partial_tv_{*,k}=0$, so the box stays full;
          \item with no diffusion, it gives $\partial_tP_*=0$;
          \item $\int R_A\,\mathrm{d}V=0$ with this estimate and with the snapshot differences, to the slot resolution. The two
                estimates agree as the resolution and the snapshot interval are refined, up to the effect of numerical
                mixing, which is reported.
        \end{enumerate}
\end{enumerate}

% =====================================================================================
\section*{References}
\begin{itemize}\small
  \item Egan, Methven, Bourne \& Cullen (2025), A definition of the background state of the atmosphere using optimal
        transport, arXiv:2504.11141.
  \item Harris \& Tailleux (2018), Assessment of algorithms for computing moist available potential energy, QJRMS 144, 1501--1510.
  \item Hieronymus \& Nycander (2015), Finding the minimum potential energy state by adiabatic parcel rearrangements with a
        nonlinear equation of state: an exact solution in polynomial time, JPO 45, 1843--1857.
  \item Hoffman (1963), On simple linear programming problems, in \emph{Convexity}, Proc.\ Symp.\ Pure Math.\ 7, AMS,
        317--327.
  \item Huang (2005), Available potential energy in the world's oceans, J.\ Mar.\ Res.\ 63, 141--158.
  \item Il{\i}cak, Adcroft, Griffies \& Hallberg (2012), Spurious dianeutral mixing and the role of momentum closure, Ocean
        Modelling 45--46, 37--58.
  \item Saenz, Tailleux, Butler, Hughes \& Oliver (2015), Estimating Lorenz's reference state in an ocean with a nonlinear
        equation of state for seawater, JPO 45, 1242--1257.
  \item Stansifer, O'Gorman \& Holt (2017), Accurate computation of moist available potential energy with the Munkres
        algorithm, QJRMS 143, 288--292.
  \item Su \& Ingersoll (2016), On the minimum potential energy state and the eddy size-constrained APE density, JPO 46,
        2663--2674.
  \item Tailleux (2013), Available potential energy density for a multicomponent Boussinesq fluid with arbitrary nonlinear
        equation of state, JFM 735, 499--518.
  \item Tailleux (2018), Local available energetics of multicomponent compressible stratified fluids, JFM 842, R1.
  \item Tseng \& Ferziger (2001), Mixing and available potential energy in stratified flows, Phys.\ Fluids 13, 1281--1293.
  \item Winters, Lombard, Riley \& D'Asaro (1995), Available potential energy and mixing in density-stratified fluids,
        JFM 289, 115--128.
  \item Winters \& D'Asaro (1996), Diascalar flux and the rate of fluid mixing, JFM 317, 179--193.
  \item Wong, Tailleux \& Gray (2016), The computation of reference state and APE production by diabatic processes in an
        idealized tropical cyclone, QJRMS 142, 2646--2657.
\end{itemize}

\end{document}
